\documentclass[11pt]{article}
\usepackage[a4paper,margin=0.85in]{geometry}
\usepackage{amsmath,amssymb,amsthm,mathtools}
\usepackage{enumitem}
\usepackage{xcolor}
\usepackage{booktabs}
\usepackage{lmodern}
\usepackage{hyperref}

\hypersetup{colorlinks=true,linkcolor=blue,urlcolor=blue}
\setlength{\parindent}{0pt}
\setlength{\parskip}{5pt}

\newcommand{\yvec}{\mathbf y}
\newcommand{\bvec}{\boldsymbol\beta}
\newcommand{\evec}{\boldsymbol\varepsilon}
\newcommand{\xmat}{\mathbf X}
\newcommand{\R}{\mathbb R}

\title{Mock Proof Questions and Solutions\\[2mm]
\large Multiple Linear Regression and Polynomial Regression}
\author{}
\date{}

\begin{document}
\maketitle

\begin{center}
\textit{Focus: ordinary derivatives, matrix transformations, gradients, and substitution.}
\end{center}

\section*{Notation and derivative rules}
For the MLR model, let
\[
\yvec=\xmat\bvec+\evec,
\qquad
\yvec\in\R^{n\times1},\quad
\xmat\in\R^{n\times p},\quad
\bvec\in\R^{p\times1}.
\]
The $i$th row of $\xmat$ is
\[
\mathbf x_i^\top=(1,x_{i1},\ldots,x_{i,p-1}).
\]

The most useful matrix identities are
\[
(AB)^\top=B^\top A^\top,
\qquad
(A+B)^\top=A^\top+B^\top,
\qquad
(A^{-1})^\top=(A^\top)^{-1},
\]
and
\[
\xmat^\top\xmat\text{ is symmetric because }
(\xmat^\top\xmat)^\top=\xmat^\top\xmat.
\]

For a vector $\bvec$ and a constant vector $a$,
\[
\nabla_{\bvec}(a^\top\bvec)=a.
\]
For a constant matrix $A$,
\[
\nabla_{\bvec}(\bvec^\top A\bvec)=(A+A^\top)\bvec.
\]
If $A$ is symmetric, this becomes
\[
\nabla_{\bvec}(\bvec^\top A\bvec)=2A\bvec.
\]

\newpage
\section*{Question 1 --- Derive the MLR OLS estimator}
Consider the model
\[
\yvec=\xmat\bvec+\evec.
\]
Define the ordinary least-squares loss
\[
L(\bvec)=(\yvec-\xmat\bvec)^\top(\yvec-\xmat\bvec).
\]
\begin{enumerate}[label=(\alph*)]
\item Expand $L(\bvec)$.
\item Compute the gradient $\nabla_{\bvec}L$ clearly.
\item Set the gradient equal to zero and derive the normal equations.
\item Assuming $\xmat^\top\xmat$ is invertible, derive the OLS estimator.
\end{enumerate}

\subsection*{Solution}
First expand the loss:
\begin{align*}
L(\bvec)
&=(\yvec^\top-\bvec^\top\xmat^\top)(\yvec-\xmat\bvec)\\
&=\yvec^\top\yvec-\yvec^\top\xmat\bvec
-\bvec^\top\xmat^\top\yvec
+\bvec^\top\xmat^\top\xmat\bvec.
\end{align*}
The two middle terms are scalars, and
\[
\yvec^\top\xmat\bvec
=\left(\yvec^\top\xmat\bvec\right)^\top
=\bvec^\top\xmat^\top\yvec.
\]
Therefore,
\[
L(\bvec)=\yvec^\top\yvec
-2\bvec^\top\xmat^\top\yvec
+\bvec^\top\xmat^\top\xmat\bvec.
\]

Now differentiate term by term:
\[
\nabla_{\bvec}(\yvec^\top\yvec)=0,
\]
\[
\nabla_{\bvec}(-2\bvec^\top\xmat^\top\yvec)
=-2\xmat^\top\yvec,
\]
and, because $\xmat^\top\xmat$ is symmetric,
\[
\nabla_{\bvec}
\left(\bvec^\top\xmat^\top\xmat\bvec\right)
=2\xmat^\top\xmat\bvec.
\]
Hence,
\[
\boxed{
\nabla_{\bvec}L
=-2\xmat^\top\yvec+2\xmat^\top\xmat\bvec
}.
\]
At the minimizer $\widehat{\bvec}$, the gradient is zero:
\[
-2\xmat^\top\yvec
+2\xmat^\top\xmat\widehat{\bvec}=0.
\]
Divide by $2$:
\[
\xmat^\top\xmat\widehat{\bvec}=\xmat^\top\yvec.
\]
These are the normal equations. Premultiplying by $(\xmat^\top\xmat)^{-1}$ gives
\[
\boxed{
\widehat{\bvec}=(\xmat^\top\xmat)^{-1}\xmat^\top\yvec
}.
\]

\textbf{Dimension check:}
\[
(p\times n)(n\times n)(n\times1)=p\times1,
\]
so the result has the correct dimension for $\widehat{\bvec}$.

\newpage
\section*{Question 2 --- Derive the MLR normal equations using ordinary derivatives}
For
\[
y_i=\beta_0+\beta_1x_{i1}+\cdots+\beta_{p-1}x_{i,p-1}+\varepsilon_i,
\]
define
\[
L(\beta_0,\ldots,\beta_{p-1})
=\sum_{i=1}^n
\left(y_i-\beta_0-\sum_{j=1}^{p-1}\beta_jx_{ij}\right)^2.
\]
Show using only ordinary partial derivatives that the normal equations are
\[
\sum_{i=1}^n x_{ij}
\left(y_i-\sum_{k=0}^{p-1}x_{ik}\beta_k\right)=0,
\qquad j=0,1,\ldots,p-1,
\]
where $x_{i0}=1$. Then show that these equations are exactly
\[
\xmat^\top\xmat\bvec=\xmat^\top\yvec.
\]

\subsection*{Solution}
Let
\[
r_i=y_i-\sum_{k=0}^{p-1}x_{ik}\beta_k.
\]
For a fixed coefficient $\beta_j$,
\[
\frac{\partial r_i}{\partial\beta_j}=-x_{ij}.
\]
Using the chain rule,
\begin{align*}
\frac{\partial L}{\partial\beta_j}
&=\sum_{i=1}^n2r_i\frac{\partial r_i}{\partial\beta_j}\\
&=-2\sum_{i=1}^n x_{ij}r_i.
\end{align*}
Set this equal to zero:
\[
\sum_{i=1}^n x_{ij}
\left(y_i-\sum_{k=0}^{p-1}x_{ik}\beta_k\right)=0.
\]
Rearrange:
\[
\sum_{k=0}^{p-1}
\left(\sum_{i=1}^n x_{ij}x_{ik}\right)\beta_k
=\sum_{i=1}^n x_{ij}y_i.
\]

The $(j,k)$ entry of $\xmat^\top\xmat$ is
\[
(\xmat^\top\xmat)_{jk}=\sum_{i=1}^n x_{ij}x_{ik},
\]
and the $j$th entry of $\xmat^\top\yvec$ is
\[
(\xmat^\top\yvec)_j=\sum_{i=1}^n x_{ij}y_i.
\]
Thus all $p$ scalar equations combine into
\[
\boxed{\xmat^\top\xmat\bvec=\xmat^\top\yvec}.
\]

\newpage
\section*{Question 3 --- Prove that the OLS estimator is unbiased}
Assume
\[
\yvec=\xmat\bvec+\evec,
\qquad E(\evec)=\mathbf0,
\]
and
\[
\widehat{\bvec}=(\xmat^\top\xmat)^{-1}\xmat^\top\yvec.
\]
Prove that
\[
E(\widehat{\bvec})=\bvec.
\]

\subsection*{Solution}
Substitute the model into the estimator:
\begin{align*}
\widehat{\bvec}
&=(\xmat^\top\xmat)^{-1}\xmat^\top
(\xmat\bvec+\evec)\\
&=(\xmat^\top\xmat)^{-1}\xmat^\top\xmat\bvec
+(\xmat^\top\xmat)^{-1}\xmat^\top\evec\\
&=\bvec+(\xmat^\top\xmat)^{-1}\xmat^\top\evec.
\end{align*}
Taking expectations and treating $\xmat$ and $\bvec$ as fixed:
\begin{align*}
E(\widehat{\bvec})
&=E\left[\bvec+(\xmat^\top\xmat)^{-1}\xmat^\top\evec\right]\\
&=\bvec+(\xmat^\top\xmat)^{-1}\xmat^\top E(\evec)\\
&=\bvec+(\xmat^\top\xmat)^{-1}\xmat^\top\mathbf0\\
&=\boxed{\bvec}.
\end{align*}
Thus OLS is unbiased: across repeated samples, the average estimator equals the true coefficient vector.

\section*{Question 4 --- Residual orthogonality}
Let
\[
\widehat{\yvec}=\xmat\widehat{\bvec},
\qquad
\widehat{\mathbf e}=\yvec-\xmat\widehat{\bvec}.
\]
Use the derivatives of the OLS loss to prove
\[
\boxed{\xmat^\top\widehat{\mathbf e}=\mathbf0}.
\]
Interpret this result for the intercept and for each predictor.

\subsection*{Solution}
The gradient is
\[
\nabla_{\bvec}L=-2\xmat^\top(\yvec-\xmat\bvec).
\]
At $\widehat{\bvec}$, the gradient is zero:
\[
-2\xmat^\top(\yvec-\xmat\widehat{\bvec})=\mathbf0.
\]
Since $\widehat{\mathbf e}=\yvec-\xmat\widehat{\bvec}$,
\[
\boxed{\xmat^\top\widehat{\mathbf e}=\mathbf0}.
\]
If the first column of $\xmat$ is a column of ones, then
\[
\sum_{i=1}^n\widehat e_i=0.
\]
For predictor $j$,
\[
\sum_{i=1}^n x_{ij}\widehat e_i=0.
\]
Thus the residuals are orthogonal to the intercept column and to every predictor column.

\newpage
\section*{Question 5 --- Quadratic polynomial regression}
Consider
\[
y_i=\beta_0+\beta_1x_i+\beta_2x_i^2+\varepsilon_i.
\]
The loss is
\[
L(\beta_0,\beta_1,\beta_2)
=\sum_{i=1}^n
\left(y_i-\beta_0-\beta_1x_i-\beta_2x_i^2\right)^2.
\]
\begin{enumerate}[label=(\alph*)]
\item Derive all three normal equations using ordinary partial derivatives.
\item Solve the first equation for $\beta_0$.
\item After substitution, express the remaining two equations using centered transformed predictors.
\end{enumerate}

\subsection*{Solution}
Let
\[
r_i=y_i-\beta_0-\beta_1x_i-\beta_2x_i^2.
\]
The three derivatives are
\[
\frac{\partial L}{\partial\beta_0}=-2\sum_i r_i,
\qquad
\frac{\partial L}{\partial\beta_1}=-2\sum_i x_ir_i,
\qquad
\frac{\partial L}{\partial\beta_2}=-2\sum_i x_i^2r_i.
\]
Setting them to zero gives
\[
\sum_i y_i-n\beta_0-\beta_1\sum_i x_i-\beta_2\sum_i x_i^2=0,
\]
\[
\sum_i x_iy_i-\beta_0\sum_i x_i-\beta_1\sum_i x_i^2-\beta_2\sum_i x_i^3=0,
\]
\[
\sum_i x_i^2y_i-\beta_0\sum_i x_i^2-\beta_1\sum_i x_i^3-\beta_2\sum_i x_i^4=0.
\]
Equivalently,
\[
\boxed{n\beta_0+\beta_1\sum_i x_i+\beta_2\sum_i x_i^2=\sum_i y_i},
\]
\[
\boxed{\beta_0\sum_i x_i+\beta_1\sum_i x_i^2+\beta_2\sum_i x_i^3=\sum_i x_iy_i},
\]
\[
\boxed{\beta_0\sum_i x_i^2+\beta_1\sum_i x_i^3+\beta_2\sum_i x_i^4=\sum_i x_i^2y_i}.
\]

From the first equation,
\[
\boxed{\beta_0=\bar y-\beta_1\bar x-\beta_2\overline{x^2}},
\qquad
\overline{x^2}=\frac1n\sum_i x_i^2.
\]

For the substitution step, define
\[
u_i=x_i-\bar x,
\qquad
v_i=x_i^2-\overline{x^2},
\qquad
w_i=y_i-\bar y.
\]
Because $\sum_i u_i=\sum_i v_i=\sum_i w_i=0$, substituting the intercept into the second and third normal equations gives
\[
\sum_i u_i^2\,\beta_1+\sum_i u_iv_i\,\beta_2=\sum_i u_iw_i,
\]
\[
\sum_i u_iv_i\,\beta_1+\sum_i v_i^2\,\beta_2=\sum_i v_iw_i.
\]
Define
\[
A_{11}=\sum_i u_i^2,\quad
A_{12}=\sum_i u_iv_i,\quad
A_{22}=\sum_i v_i^2,
\]
\[
B_1=\sum_i u_iw_i,\qquad B_2=\sum_i v_iw_i.
\]
Then
\[
A_{11}\beta_1+A_{12}\beta_2=B_1,
\qquad
A_{12}\beta_1+A_{22}\beta_2=B_2.
\]
If $D=A_{11}A_{22}-A_{12}^2\neq0$, then
\[
\boxed{\widehat\beta_1=\frac{B_1A_{22}-B_2A_{12}}{D}},
\qquad
\boxed{\widehat\beta_2=\frac{B_2A_{11}-B_1A_{12}}{D}},
\]
and
\[
\boxed{\widehat\beta_0=\bar y-\widehat\beta_1\bar x-\widehat\beta_2\overline{x^2}}.
\]

\newpage
\section*{Question 6 --- Polynomial regression as MLR}
For the degree-$d$ polynomial model
\[
y_i=\beta_0+\beta_1x_i+\beta_2x_i^2+\cdots+\beta_dx_i^d+\varepsilon_i,
\]
\begin{enumerate}[label=(\alph*)]
\item Construct the design matrix $\xmat$.
\item Show that the model can be written as $\yvec=\xmat\bvec+\evec$.
\item Derive the OLS estimator from the matrix loss.
\end{enumerate}

\subsection*{Solution}
The coefficient vector is
\[
\bvec=(\beta_0,\beta_1,\ldots,\beta_d)^\top.
\]
The design matrix is
\[
\xmat=
\begin{pmatrix}
1&x_1&x_1^2&\cdots&x_1^d\\
1&x_2&x_2^2&\cdots&x_2^d\\
\vdots&\vdots&\vdots&&\vdots\\
1&x_n&x_n^2&\cdots&x_n^d
\end{pmatrix}.
\]
Its $i$th row times $\bvec$ is
\[
\begin{pmatrix}1&x_i&x_i^2&\cdots&x_i^d\end{pmatrix}
\begin{pmatrix}\beta_0\\\beta_1\\\beta_2\\\vdots\\\beta_d\end{pmatrix}
=\beta_0+\beta_1x_i+\cdots+\beta_dx_i^d.
\]
Thus the polynomial model is an MLR model with transformed predictors
\[
x_i,\;x_i^2,\;\ldots,\;x_i^d.
\]
Using
\[
L(\bvec)=(\yvec-\xmat\bvec)^\top(\yvec-\xmat\bvec),
\]
the same expansion and gradient calculation as in Question 1 gives
\[
\nabla_{\bvec}L
=-2\xmat^\top\yvec+2\xmat^\top\xmat\bvec.
\]
Set this equal to zero:
\[
\xmat^\top\xmat\widehat{\bvec}=\xmat^\top\yvec.
\]
Therefore,
\[
\boxed{
\widehat{\bvec}=(\xmat^\top\xmat)^{-1}\xmat^\top\yvec
}.
\]
Polynomial regression is nonlinear in $x$, but it is linear in the unknown parameters $\beta_0,\ldots,\beta_d$.

\section*{Exam checklist}
\begin{enumerate}[leftmargin=*]
\item Define the residual clearly.
\item Differentiate with respect to each unknown coefficient.
\item Set each derivative equal to zero.
\item For matrix proofs, remember $(AB)^\top=B^\top A^\top$.
\item Combine scalar equations into $\xmat^\top\xmat\widehat{\bvec}=\xmat^\top\yvec$.
\item Solve the intercept equation first when substitution is requested.
\item Check dimensions before multiplying matrices.
\item In polynomial regression, use the transformed columns $1,x,x^2,\ldots,x^d$.
\end{enumerate}

\end{document}

